
For the definition of the general case's procedure, we need to introduce some notation and definitions. We assume that we are given some input data, which consists of the lower triangular bidding matrix $\Set{b_{i,j} \mid i \geq j \in \{1, 2, ..., n\}}$, and some advice variables $\Gamma$. First, we will define the following quantities related to the bidding matrix:

\begin{definition}
    For each client $j\in\{1,\dots,n\}$, define its \emph{first nonzero bidding time}
    \begin{equation*}
        f_j := \min\{i\in\{j,\dots,n\} : b_{i,j}=1\}\,.
    \end{equation*}
    This quantity is well defined by Observation~\ref{obs:b1}. For each $i\in\{1,\dots,n\}$, we also write
    \begin{equation*}
        \ones{i} := \{j : f_j=i\}\,.
    \end{equation*}
    Thus $\ones{i}$ is the set of clients whose first non-zero bid appears at time $i$, and by lower triangularity we have $\ones{i}\subseteq\{1,\dots,i\}$.
\end{definition}

As for the advice, let us define it.

\begin{definition}[Advice Variables]\label{def:auxiliary-input}
    For a fixed bidding matrix $\{b_{i,j}\mid i\ge j\}$, an \emph{advice} for the procedure consists of
    \begin{equation*}
        t_3^\star \in \{1,\dots,n\}\,, \qquad
        \{\mu_j^\star\}_{j\in\ones{t_3^\star}}\,,
    \end{equation*}
    and, in the special case $t_3^\star=n$, an additional scalar
    \begin{equation*}
        \hg_n^\star \in [0,1]\,.
    \end{equation*}
    % The index $t_3^\star$ specifies the frontier at which the construction switches from the trivial Phase 3 assignments to the recursive construction of the $\hg_i$'s. The values $\mu_j^\star$ prescribe the initial $D$-membership of the clients on that frontier, namely those $j$ with $f_j=t_3^\star$. When $t_3^\star=n$, the value $\hg_n^\star$ is used as the initial boundary value for $\hg_n$; if $t_3^\star<n$, no additional $\hg$-input is needed.
\end{definition}


The $\{\mu_j^\star\}_{j\in\ones{t_3^\star}}$ must satisfy the following properties, necessary for the construction to work and which will be used in the proof of correctness in Section~\ref{sec:general-case-proof}.

\begin{property}[Admissible frontier input]\label{property:admissible-frontier-input}
    Fix $t_3^\star\in\{1,\dots,n\}$ and write $\ones{t_3^\star}=\{j_1<\dots<j_k\}$. We call a family $\{\mu_j^\star\}_{j\in\ones{t_3^\star}}$ \emph{admissible} if it is non-increasing in $j$, has at most one fractional entry, and, when $t_3^\star=n$, satisfies $\mu_j^\star=1$ for all $j\in\ones{n}$.

    Equivalently, there exist $r\in\{0,\dots,k\}$ and $\lambda\in[0,1]$ such that
    \begin{equation*}
        \mu_{j_p}^\star=
        \begin{cases}
            1, & p\le r\,,\\
            \lambda, & p=r+1\le k\,,\\
            0, & p>r+1\,.
        \end{cases}
    \end{equation*}
    where $r=k$ in the case of $t_3^\star=n$. 
\end{property}

Throughout the procedure, we will separate and treat the clients $j$ differently according to their position with respect to two thresholds $t_1$ and $t_2$, generated by the procedure itself, and the input one $t_3^\star$. Putting it very simply, each client $j$ will belong to one of the following four groups:
\begin{itemize}
    \item Group $A$: clients with $j < t_1$ and $f_j< t_3^\star$.
    \item Group $B$: clients with $t_1 < j \le t_2$ and $f_j< t_3^\star$.
    \item Group $C$: clients with $t_2 < j \le t_3^\star$ and $f_j< t_3^\star$.
    \item Group $D$: clients with $f_j > t_3^\star$.
\end{itemize}
Notice that we did not specify what happens to clients with $f_j = t_3^\star$ or to the $j=t_1$ client; their behaviours are more delicate, and they might belong to more than one group at the same time. To work with this, we will use the following notation.
\begin{definition}[Membership weights]
    For each client $j\in\{1,\dots,n\}$, the construction produces four auxiliary scalars
    \begin{equation*}
        \mu_j^A,\mu_j^B,\mu_j^C,\mu_j^D\in[0,1]\,,
    \end{equation*}
    called the \emph{membership weights} of client $j$ in the classes $A,B,C,D$. For every client $j$, they satisfy
    \begin{equation*}
        \mu_j^A+\mu_j^B+\mu_j^C+\mu_j^D=1\,.
    \end{equation*}
\end{definition}
For the $t_1$ client, this construction will be delicate, but for the $f_j=t_3^\star$ clients, we will set $\mu_j^D=\mu_j^\star$, assigning the rest of the mass to $A$ (if $j<t_1$), to $B$ (if $t_1<j\le t_2$), and to $C$ (if $t_2<j\le t_3^\star$). In particular, if $t_3^\star=n$, we will have $\mu_j^D=1$ for all clients with $f_j=t_3^\star$, meaning that all those clients belong to group $D$.

Through these inputs, the procedure described in Algorithm~\ref{alg:general-case-procedure-hgi} generates a feasible solution for the dual of the factor-revealing LP using the function $\beta_{\Mcoeff}(\alpha,d) = [\alpha - \Mcoeff d]^+$, by first generating the $\hg_i$ values in Algorithm~\ref{alg:general-case-procedure-hgi} and then using them to generate the $g_{i,j}, s_{i,j}, t_{i,j}$ values in Table~\ref{tab:general-case-procedure}. If the advice is not correct, the procedure can fail and return an empty output ($\perp$). 

For readability, we will use the following definitions:
\begin{definition}
    We define the following shorthand variables which will help in our analysis, all a function of the value of $\Mcoeff$:
    \begin{align*}
        \Mcoeffmone \coloneqq \Mcoeff - 1\,,
        \qquad
        \MRttwo \coloneqq \Mcoeffmone + \Mcoeffmone^2\,,
        \qquad
        \Mtone \coloneqq 2\cdot \Mcoeffmone + \Mcoeffmone^2\,,
        \qquad
        \MRtone \coloneqq 1-\Mtone
    \end{align*}
\end{definition}
from where the following properties follow:
\begin{property}\label{property:coefficients}
    For any values of $\Mcoeff$, the following properties hold:
    \begin{equation*}
        \Mcoeff \Mcoeffmone = \MRttwo, \qquad \Mcoeff + 1 - \MRttwo = \Mrho, \qquad \Mfrac \leq 1\,.
    \end{equation*}
\end{property}
Thanks to the specific choice of $\Mcoeff$ we made, we also have the following property on the ordering of the variables:
\begin{property}\label{property:thresholds}
    For $\Mcoeff \in [1,1.224]$, we have $\Mtone \in [0,0.5)$ and thus $\Mcoeffmone \leq \MRttwo \leq \Mtone < 0.5 < \MRtone$.
\end{property}
Finally, we note the following technical property that the variables satisfy:
\begin{property}\label{property:function-tau}
    For $\Mcoeff \in [1,1.224]$, we have that for $x\in[0,1-\Mtone]$, the following holds:
    \begin{equation*}    
        \frac{1-\Mtone-x}{\Mtone + x} \geq \left[1-\frac{x}{\MRttwo}\right]^+\,.
    \end{equation*}
\end{property}
% \toni{might have to write a proof for this in appendix.}

This last technical property is the bottleneck for the choice of $\Mcoeff$, as it is necessary for our arguments to work, and thus it is the bottleneck for the approximation factor $\Mrho$ that we can achieve through our technique.

We will also use the prefix sums to express:
\begin{equation*}
    R_i \coloneqq \sum_{\ell=i+1}^n \hg_\ell, \qquad\text{and}\qquad R_i^{(2)} \coloneqq \sum_{\ell=i+1}^n \hg_\ell^{(2)}\,.
\end{equation*}
In particular, $R_n=R_n^{(2)}=0$, and $R_{i-1}=R_i+\hg_i$, $R_{i-1}^{(2)}=R_i^{(2)}+\hg_i^{(2)}$.

\begin{algorithm}[H]
\caption{Construction of the profile $\hg$}\label{alg:general-case-procedure-hgi}
\begin{algorithmic}[1]
    \Require $\{b_{i,j} : i\ge j\in\{1,\dots,n\}\}$, an index $t_3^\star\in\{1,\dots,n\}$, an admissible frontier input $\{\mu_j^\star : j\in\ones{t_3^\star}\}$, and $\hg_n^\star\in[0,1]$ when $t_3^\star=n$

    \State Set
    $\mu_j^D=
    \begin{cases}
        1, & f_j>t_3^\star\,,\\
        \mu_j^\star, & j\in\ones{t_3^\star}\,,\\
        0, & f_j<t_3^\star\,.
    \end{cases}$

    \State Set $R_{t_3^\star}^{(2)}\gets 0$ and $\hg_i^{(2)}\gets R_i^{(2)}\gets 0$ for every $i>t_3^\star$

    \For{each $t_3^\star \ge i \ge 1$} \label{line:phase-2-loop}
        \State $\hg_i^{(2)}\gets
        \begin{cases}
            \hg_n^\star, & \text{if $i=t_3^\star=n$}\,,\\
            (1-\mu_i^D)\dfrac{R_i^{(2)}+\Mtone}{n-i}\,, & \text{otherwise}
        \end{cases}$
        \State $R_{i-1}^{(2)}\gets R_i^{(2)}+\hg_i^{(2)}$.
    \EndFor

    \State $t_1\gets \max\{i : R_{i-1}^{(2)}\ge \MRtone\}$ \label{line:R_i-1-smaller-0.79-condition}
    \State Set $(\hg_i,R_i)\gets (\hg_i^{(2)},R_i^{(2)})$ for every $i>t_1$ (or all $i$ if undefined) and $R_{t_1}\gets R_{t_1}^{(2)}$ (if defined).
    \If{$t_1$ does not exist} \label{line:t_1-perp-condition}
        \State \Return $\perp$
    \EndIf

    \State $t_2\gets \max\{i : R_{i-1}^{(2)}\ge \MRttwo\}$
    \State $z\gets \sum_{\ell=t_2+1}^{n}\left[1-\dfrac{R_{\ell-1}^{(2)}}{\MRttwo}\right]^+\mu_\ell^D$

    \State Define $\alpha$, $\hg_{t_1}^A$, $\hg_{t_1}^B$, $\hg_{t_1}$, and $R_{t_1-1}$ according to Definition~\ref{def:t1-boundary-split}

    \For{each $t_1 > i \geq 1$}
        \State $\hg_i \gets (1-\mu_i^D)\cdot\dfrac{1-[\MRttwo-R_{f_i}]^+}{n-i-\Mfrac\cdot z}$
        \State $R_{i-1} \gets R_i + \hg_i$
    \EndFor

    \If{$R_0 \neq 1$}
        \State \Return $\perp$
    \EndIf
\end{algorithmic}
\end{algorithm}

\begin{definition}[Boundary Split at $t_1$]\label{def:t1-boundary-split}
Assume that $t_1$ has been defined by Algorithm~\ref{alg:general-case-procedure-hgi}.
If $t_1 < n$, define
\begin{equation*}
    \alpha \coloneqq \frac{\MRtone - R_{t_1}^{(2)}}{\hg_{t_1}^{(2)}}\,,
    \qquad
    \hg_{t_1}^B \coloneqq \MRtone - R_{t_1}^{(2)}\,,
    \qquad
    \hg_{t_1}^A \coloneqq
    (1-\alpha)(1-\mu_{t_1}^D)\cdot
    \frac{1-[\MRttwo-R_{f_{t_1}}]^+}
         {n-t_1+\alpha(1-\mu_{t_1}^D)-\Mfrac\cdot z}\,.
\end{equation*}
If $t_1=n$, define instead
\begin{equation*}
    \alpha \coloneqq 1\,,
    \qquad
    \hg_n^A \coloneqq 0\,,
    \qquad
    \hg_n^B \coloneqq \hg_n^{(2)}=\hg_n^\star\,.
\end{equation*}
In both cases, we finally set
\begin{equation*}
    \hg_{t_1} \coloneqq \hg_{t_1}^A+\hg_{t_1}^B\,,
    \qquad
    R_{t_1-1} \coloneqq R_{t_1}+\hg_{t_1}\,.
\end{equation*}
\end{definition}
It is important to note the following about $\alpha$:
\begin{property}
We have $0 \leq \alpha \leq 1$.
\end{property}
% \begin{proof}
% If $t_1 = n$, then $\alpha = 1$.
% Else, we have $R_{t_1 - 1}^{(2)} = R_{t_1}^{(2)} + \hg_{t_1}^{(2)}$ together with $R_{t_1 - 1}^{(2)} \geq \MRtone$ and $R_{t_1}^{(2)} < \MRtone$. Hence $0 \leq \MRtone - R_{t_1}^{(2)} \leq \hg_{t_1}^{(2)}$. Therefore $\alpha = \frac{\MRtone - R_{t_1}^{(2)}}{\hg_{t_1}^{(2)}}$ satisfies $0 \leq \alpha \leq 1$.
% \end{proof}

With these values, we set the memberships as follows:
\begin{equation}\label{eq:membership-assignment}
    \begin{aligned}
        &\mu_j^A=
        \begin{cases}
            1-\mu_j^D, & j<t_1\,,\\
            (1-\alpha)(1-\mu_{t_1}^D), & j=t_1\,,\\
            0, & j>t_1\,,
        \end{cases}
        \qquad
        &&\mu_j^B=
        \begin{cases}
            \alpha(1-\mu_{t_1}^D), & j=t_1\,,\\
            1-\mu_j^D, & t_1<j\le t_2\,,\\
            0, & \text{otherwise}\,,
        \end{cases}
        \\
        &\mu_j^C=
        \begin{cases}
            1-\mu_j^D, & t_2<j\le t_3^\star\,,\\
            0, & \text{otherwise}\,,
        \end{cases} 
        \qquad
        &&\mu_j^D=
        \begin{cases}
            1, & f_j>t_3^\star,\\
            \mu_j^\star, & j\in\ones{t_3^\star}\,,\\
            0, & f_j<t_3^\star\,.
        \end{cases}
    \end{aligned}
\end{equation}
In particular, the only non-integral memberships are the boundary input value on $\ones{t_3^\star}$ and the split of the client $t_1$ between $A$ and $B$.

Once $\hat g$ and $R$ are fixed, we abbreviate
\begin{equation*}
    \chi_j := \left[1-\frac{R_{j-1}}{\theta}\right]^+\,.
\end{equation*}

The matrices $g_{i,j}$, $s_{i,j}$, and $t_{i,j}$ are then assigned column-wise according to Table~\ref{tab:general-case-procedure}. Every entry not listed there is set to $0$. The variables $\delta_{i,j}$ and $\sigma_{i,j}$ are auxiliary variables that need to satisfy Property~\ref{property:delta} and Property~\ref{property:sigma}, respectively, and whose existence will be proved in the proof of Lemma~\ref{lem:general-case-procedure}.

\begin{table}[H]
\centering
\small
\caption{Column-wise assignment of $(g_{i,j},s_{i,j},t_{i,j})$ once $\hat g$ and $R$ are fixed.}
\label{tab:general-case-procedure}
\renewcommand{\arraystretch}{1.25}
\begin{tabularx}{\textwidth}{@{}>{\RaggedRight\arraybackslash}p{0.14\textwidth}>{\RaggedRight\arraybackslash}p{0.27\textwidth}>{\RaggedRight\arraybackslash}p{0.28\textwidth}>{\RaggedRight\arraybackslash}X@{}}
\toprule
Group Membership & Generated credit $g_{i,j}$ & Pool transfer $s_{i,j}$ & $t$-entry and extra update \\
\midrule
$j\in A$ &
$g_{i,j}^A = \hg_i-\delta_{i,j}$ &
$s_{i,j}^A = \hg_i-\sigma_{i,j}$ for $i\ge j$ &
$t_{f_j,j}^A = 0$ \\
\addlinespace[0.4em]
$j\in B$ &
$g_{i,j}^B = \hg_i$ &
\begin{tabular}{@{}l@{}}
    $s_{i,j}^B = \hg_i$ for $i>j$ \\
    $s_{f_j+1,j}^B=\hg_{f_j+1} + \Mcoeffmone$
\end{tabular}
& $t_{f_{j},j}^B= [1-R_{j-1} - \Mtone]^+$
\\
\addlinespace[0.4em]
$j\in C$ &
$g_{i,j}^C = \hg_i$ &
\begin{tabular}{@{}l@{}}
    $s_{i,j}^C = \hg_i$ for $i>j$ \\
    $s_{f_j+1,j}^C=\hg_{f_j+1} + \Mcoeffmone$
\end{tabular}
&$t_{f_{j},j}^C=1-R_{j-1} - \Mtone$
\\
\addlinespace[0.4em]
$j\in D$ &
$g_{i,j}^D = \begin{cases}
    \hg_i - \chi_j\Mfrac\hg_i, & i<t_1,\\
    \hg_i - \chi_j\Mfrac\hg_{t_1}^A, & i=t_1,\\
    \hg_i, & i>t_1
\end{cases}$ &
$s_{i,j}^D = 0$ for $i>j$ &
$t_{f_{j},j}^D=1-R_{j-1}$
\\
\bottomrule
\end{tabularx}
\end{table}

For the cases that the client is fractionally assigned to more than one group, the generated $g_{i,j}$, $s_{i,j}$, and $t_{i,j}$ values are the corresponding convex combinations of the values generated by each group:
\begin{equation}\label{formula:fractional-assignments-convex-combination}
    \begin{aligned}
        g_{i,j} &= \mu_j^A g_{i,j}^A + \mu_j^B g_{i,j}^B + \mu_j^C g_{i,j}^C + \mu_j^D g_{i,j}^D\,,\\
        s_{i,j} &= \mu_j^A s_{i,j}^A + \mu_j^B s_{i,j}^B + \mu_j^C s_{i,j}^C + \mu_j^D s_{i,j}^D\,,\\
        t_{i,j} &= \mu_j^A t_{i,j}^A + \mu_j^B t_{i,j}^B + \mu_j^C t_{i,j}^C + \mu_j^D t_{i,j}^D\,.
    \end{aligned}
\end{equation}

As stated before, there are two variables $\delta_{i,j}$ and $\sigma_{i,j}$ that need to satisfy some properties:

\begin{property}\label{property:delta}
    The values of $\delta_{i,j}$ for $\mu_j^A>0$ need to satisfy the following properties:
    \begin{itemize}
        \item $\forall \ i, j$: $0 \le \delta_{i,j} \le \hg_i$.
        \item For $i\leq f_j$: $\delta_{i,j}=0$.
        \item $\sum_{i=f_j+1}^{t_3^\star} \delta_{i,j} = \Mcoeffmone \cdot \left[1 -\frac{[\MRttwo-R_{f_j}]^+}{\MRttwo}\right]$.
    \end{itemize}
\end{property}

\begin{property}\label{property:sigma}
    The values of $\sigma_{i,j}$ for $\mu_j^A>0$ need to satisfy the following properties:
    \begin{itemize}
        \item $\forall \ i, j$: $0 \le \sigma_{i,j} \le \hg_i$.
        \item For $i < j$ or $i > f_j$: $\sigma_{i,j}=0$.
        \item $\sum_{i=j}^{f_j} \sigma_{i,j} = [\MRttwo-R_{f_j}]^+$.
    \end{itemize}
\end{property}
